\chapter{Practice Problem Solutions}
\label{app:answers}

This appendix provides answers to the text exercises. Please try it yourself and then check.

\section{Chapter 1 Answers}

\textbf{1-1.}Rate ChatGPT’s answer.

Best answer: ``Language models are trained to predict next tokens, but with sufficiently large models and diverse data, they implicitly learn question-answer patterns. When the user inputs a question in the form of a question, the model generates the answer pattern that follows the question in the training data.''

\textbf{1-2.}``The sun sets in the east...'' Next word probability.

For humans, ``soaring'' is rated at over 99\%. Reasons: (1) everyday observation, (2) linguistic collocation (``rising from the east'' is a fixed expression), (3) strong constraints of context.

\textbf{1-3.}N-gram sparsity problem.

For N-gram to predict the next word ``polar bear dance'', the next word ``polar bear dance'' must have appeared in the training data. If this 3-gram has not been learned, the probability becomes 0, and general words such as ``more'' are recommended through smoothing.

\textbf{1-4.}Why transformers are better than RNNs for parallel processing.

RNN requires sequential processing because the calculation of point$t$depends on the hidden state of point$t-1$. Transformers process all positions simultaneously with attention, making them optimal for GPU parallelization.

\textbf{1-5.}Chinchilla reference data scale.

If the parameters are increased by 10 times, the data must also be increased by about 10 times (maintaining a 20 token/parameter ratio). In other words, a 10x model requires 200x data.

\section{Chapter 3 Answers}

\textbf{3-1.}What if we merge less frequent pairs?

Patterns that do not appear frequently become tokens, making the vocabulary inefficient. When rare pairs like ``xz'' are merged, they result in tokens that are rarely used in actual text.

\textbf{3-2.}Unigram penalty changed to\code{-5.0}.

Since the penalty for substrings not in the vocabulary is small, Viterbi attempts segmentation more actively. As a result, the tendency to break down into many smaller tokens.

\textbf{3-3.}Korean BPE, vocabulary 100 vs 1000.

Vocabulary 100: Break down ``hello'' into tokens of 5 or more bytes. Very inefficient.
Vocabulary 1000: Able to learn subwords such as ``hello'' and ``hello''. Significantly shortened sequence length.

\textbf{3-4.}Viterbi length reduced to upper limit of 4.

Tokens longer than 5 characters disappear. Words like ``hello'' are broken down into ``he'', ``ll'', and ``o'', increasing the sequence length.

The need for\textbf{3-5.}\code{<bos>}.

The model will work without it, but if you explicitly mark the beginning of a sequence, the model will recognize that ``a new sequence starts from here''. Also useful for distinguishing between multiple sequences during batch processing.

\section{Chapter 4 Answers}

\textbf{4-1.}When the embedding dimension is too small or large.

If it is small (8): It is difficult to distinguish between words, making it impossible to learn semantic similarity.
Larger (4096): risk of parameter explosion, running out of memory, and overfitting. Usually 256$\sim$4096 used.

\textbf{4-2.}Sinusoidal constant 100$\to$100.

As the cycle becomes shorter, it becomes difficult to distinguish even close locations. Lowered the ``resolution'' of location encoding.

\textbf{4-3.}Scaling omitted.

Location embeddings dominate token embeddings, causing the model to rely solely on location. Loss of semantic information.

\textbf{4-4.}Learnable location embeddings, max\_len exceeded.

An error occurred because the index was out of range. A method that allows extrapolation, such as RoPE, is needed.

\textbf{4-5.}Change LayerNorm$\epsilon$value.

If$\epsilon$is large, normalization becomes weak and output variance increases. If it is too small, the denominator will be close to 0, causing numerical instability.

\section{Chapter 5 Answers}

\textbf{5-1.}softmax when scaling is omitted.

When$d\_k=64$, the inner product value can go up to$\pm 64$, so softmax is saturated in one-hot. Gradient 0, learning stops.

\textbf{5-2.}Causal mask not used.

During training: the model ``sees'' future tokens and answers correctly, resulting in lower learning loss but no generalization.
During inference: There are no future tokens, so a learning-inference gap occurs. Performance plummets.

\textbf{5-3.}Number of heads 1 (single head).

Learn only one attention pattern. Decreased performance because grammar, semantics, and discourse relationships must all be processed by one head.

\textbf{5-4.}Number of heads =$d\_{model}$($d\_{head}=1$).

Each head processes only one dimension. The meaning of attention disappears and becomes a simple weighted sum. Practically the same as a linear transformation.

\textbf{5-5.}Attention matrix memory,$L=100,000$.

$100000^2 \times 2$bytes (FP16) = 20 GB. Based on single head. 32 heads = 640 GB. Virtually impossible.

\section{Chapter 6 Answers}

\textbf{6-1.} $d\_{ff} = d\_{model}$.

FFN's expansion-contraction pattern disappears, reducing non-linear expressiveness. Similar to a deeper linear transformation.

\textbf{6-2.}12th floor with no residuals.

When going backwards, learning is not possible because the gradient does not reach the previous layer. Losses are not reduced.

\textbf{6-3.}Necessity of\code{ln\_f}in Pre-LN.

Pre-LN normalizes each block input, but leaves the last block output unnormalized. Normalization is required before inserting into the LM Head.

\textbf{6-4.}SiLU instead of GELU.

Minor performance difference. SiLU ($x \cdot \sigma(x)$) also has a smooth transition. Some models (FFN in LLaMA is SwiGLU but uses SiLU as activation) are adopted.

\textbf{6-5.}100 layers without residual scaling.

The initial activation value distribution accumulates for each layer and explodes. NaN is likely to occur in the early stages of training.

\section{Chapter 7 Answers}

\textbf{7-1.}No weighted ties used, GPT-2 small.

Add$V \times d = 50257 \times 768 \approx 38.6M$parameter. Total 124M$\to$163M.

\textbf{7-2.} Warmup 0.

The gradient at the beginning of training is unstable, so there is a possibility of divergence. Especially serious in Post-LN. Pre-LN is less sensitive but still recommended.

\textbf{7-3.} Label smoothing 0.3.

Strong smoothing makes the model less confident about the correct answer. Generalization improves, but learning loss remains high. Usually 0.1 is appropriate.

\textbf{7-4.}No gradient clipping.

When a gradient explodes, the weights change significantly and the model collapses. Especially dangerous at large learning rates.

\textbf{7-5.}Is PPL 1.0 possible?

Theoretically impossible. Natural language has inherent uncertainty, so the next token cannot be 100% predicted. PPL$\approx$1 is possible in certain sentences only if the learning data is completely memorized, but it is not common.

\section{Chapter 8 Answers}

\textbf{8-1.}Causes Greedy repeating loops.

If a loop is formed where the distribution predicts ``to buy'' highly after ``store'' and ``some store'' after ``to buy'' highly, greedy cannot get out of it. This is because only the path with the highest probability is followed.

\textbf{8-2.} Temperature $\to$ 0.

If$T=0$, softmax becomes argmax and is the same as greedy.  If you are$T=0.01$, you are almost greedy. The distribution is so sharp that the diversity is 0.

\textbf{8-3.}Top-K and then Top-P.

After reducing the candidates to K using Top-K, select up to the cumulative probability P among them. More powerful filtering. Quality improves but variety decreases.

\textbf{8-4.}Sampling suitable for chatbots.

Top-P (P=0.9, T=0.7$\sim$0.9). Balance between variety and consistency. Greedy is monotonous, pure sampling lacks consistency.

\textbf{8-5.}Why speculative decoding is fast.

If the small model quickly generates K tokens, the large model verifies the K tokens at once (parallel processing). If correct, K are adopted in one step, so the number of steps in the large model is reduced to 1/K.
